\chapter{Personal notes}
\label{ch:personal}

This chapter collects ideas worked out with the Feynman technique: an idea is explained in plain
words as if to someone new to it, the places where the explanation breaks down are found, and the
explanation is repaired until it holds. Unlike the rest of these notes, each entry records the
explanation \emph{in the author's own words}, as it was said, with only typing errors corrected.
It is kept that way on purpose, because an explanation in one's own words is the one that is
easiest to recall later.

Each entry has the same parts:
\begin{description}[style=nextline]
  \item[In my own words] the explanation as given, in the violet box;
  \item[Where it broke down] the questions or steps the first explanation could not answer;
  \item[Repaired] the corrected or completed explanation, again in the author's own words;
  \item[Links] the sections of these notes and the files the idea belongs to.
\end{description}
Entries are listed newest first.

% ------------------------------------------------------------------- template
% Copy this block for every new entry, newest first, below this comment.
%
% \section{<Topic> \normalfont\small(<date>)}
% \label{sec:pn-<short-name>}
%
% \begin{inmywords}
% <the explanation exactly as given>
% \end{inmywords}
%
% \paragraph{Where it broke down.}
% \begin{itemize}
%   \item <question or step that did not hold>
% \end{itemize}
%
% \begin{inmywords}[title=Repaired]
% <the corrected explanation, in the author's words>
% \end{inmywords}
%
% \paragraph{Links.} \Cref{<label>}, \file{<file>}.
% ---------------------------------------------------------------------------

\section*{No entries yet}
The first entry is added after the first idea worked out this way.
